\scanpage{98}

On aura
\[(105)\qquad\mathcal V_J^Z=\bigcup_{j\in J}\mathcal V_j^Z=\underset{d\in\mathrm{Drap}(J)}{\lim_{\longrightarrow}}\mathcal V_d^Z\]
\margin{\diagram{\begin{tikzcd}[row sep=small]\mathcal V_d^Z\arrow[d,hook,"(h)"]\\(\mathcal V_d)\end{tikzcd}} si $Z\supset X_J^*$}
\deleted{\unsure{et}}

On pose, pour deux bandes $J$ et $L$
\[
(106)\qquad\left\{\begin{aligned}
\mathcal V_{J,L}&=\mathcal V_J\cap\mathcal V_L,\qquad\mathcal V_{J,L}^Z=\mathcal V_{J,L}\cap Z\\
\mathcal V_{J\deleted{,L}}^L&=\underbrace{\mathcal V_J\cap X_L^*}_{\substack{\Vert\\\mathcal V_J^{X_L^*}\text{ avec}\\\text{notation }(102)}}\subset\mathcal V_{J,L}\cap X_L^*=\mathcal V_{J,L}^{\deleted{\unsure{Z}}}
\end{aligned}\right.
\]
\margin{avec $Z=X_L^*$}
\margin{Suggestion: $\mathcal V_{J,L}^{\unsure{\flat|\natural}}$ au lieu de $\mathcal V_J^L$\\
$\mathcal V_{i_0\ldots i_n}^{\unsure{\flat|\natural}}=\mathcal V_{i_0\ldots i_n}^{X_{i_n}^*}=\mathcal V_{i_0}\cap\cdots\cap\mathcal V_{i_n}\cap X_{i_n}^*$}

On aura donc
\[
(107)\qquad\begin{aligned}
\mathcal V_{J,L}&=\bigcup_{\substack{j\in J\\l\in L}}\underbrace{\mathcal V_j\cap\mathcal V_l}_{\mathcal V_{j,l}}\deleted{{}=\bigcup_{\substack{j\in J\\l\in L\\j\leq l}}}\\
&=\bigcup_{i\in J\cap L}\mathcal V_i\ \cup\bigcup_{\substack{j<l\\j\in J,\,l\in L}}\mathcal V_{jl}\ \cup\bigcup_{\substack{l<j\\j\in J,\,l\in L}}\mathcal V_{lj}
\end{aligned}
\]
\margin{$=\underset{\substack{d\in\mathrm{Drap}(J\cup L)\\d\cap J\ne\varnothing,\,d\cap L\ne\varnothing}}{\lim_{\longrightarrow}}\mathcal V_d$}
\deleted{\[=\underset{\unsure{d\in\mathrm{Drap}_{\geq2}(J\cup L)}}{\lim_{\longrightarrow}}\mathcal V_d\]}
donc si $J\prec L$, et $J$ et $L$ des bandes
\[
(108)\qquad\begin{aligned}
\mathcal V_{J,L}&=\bigcup_{i\in J\cap L}\mathcal V_i\ \cup\bigcup_{\substack{j<l\\j\in J,\,l\in L}}\mathcal V_{jl}\\
&=\underset{\substack{d\in\mathrm{Drap}(J\cup L)\\d\cap J\ne\varnothing,\,d\cap L\ne\varnothing}}{\lim_{\longrightarrow}}\mathcal V_d
\end{aligned}
\]
\margin{\uline{NB} ici les drapeaux en question les $i_0<\cdots<i_n$, les $i_\alpha\in J\cup L$, $i_0\in J$, $i_n\in L$, donc les $i_\alpha$ qui sont $\in J$ forment un segm. initial, les $i_\beta\in L$ un segm. final$^*$}
\margin{$^*$ La réunion des deux segm. est $\{i_0,\ldots,i_n\}$; l'intersection est \deleted{\unsure{la}} $\varnothing$ si $J\cap L=\varnothing$\ldots}
\margin{97}
